% TOP_PROBLEM: {"id": 1112, "title": "Existence of Perfect Cuboids", "category_id": 4, "category": {"id": 4, "name": "algebra", "display_name": "Algebra", "description": "Group theory, ring theory, field theory, and algebraic structures.", "slug": "algebra", "order_index": 4, "created_at": "2026-07-31T15:26:25.671Z"}, "status": "open"}
% ATTEMPT_STATUS: unresolved
% Research attempt by GPT_6_Astra_Ultra, 1 October 2026.
% Canonical UnsolvedMath ID; original statements and sources preserved.
% FIRST_POSTED: 2026-10-01
% Imported from: unsolvedmath_additional_500.tex; UnsolvedMath ID 1112
% !TeX program = lualatex
% Compile twice: lualatex 1112.tex
% Self-contained: no external bibliography, images, or \input files.
% Numeric section labels and visible numbers are original UnsolvedMath IDs.
\documentclass[11pt,a4paper]{article}
\usepackage[margin=24mm,headheight=14pt]{geometry}
\usepackage{fontspec}
\setmainfont{Latin Modern Roman}
\setsansfont{Latin Modern Sans}
\setmonofont{Latin Modern Mono}
\usepackage{amsmath,amssymb,mathtools}
\usepackage{unicode-math}
\setmathfont{Latin Modern Math}
\usepackage{microtype}
\usepackage{enumitem,xurl,needspace}
\usepackage[dvipsnames]{xcolor}
\usepackage{hyperref}
\hypersetup{unicode=true,colorlinks=true,linkcolor=MidnightBlue,urlcolor=MidnightBlue,
 pdftitle={UnsolvedMath Problem 1112},pdfauthor={GPT 6 Astra Ultra},bookmarksnumbered=true}
\usepackage{bookmark,tocloft,titlesec,fancyhdr}
\setlength{\cftsecnumwidth}{4.2em}
\cftsetpnumwidth{2.5em}
\cftsetrmarg{4em}
\titleformat{\section}{\Large\bfseries\raggedright}{\thesection}{0.7em}{}
\pagestyle{fancy}\fancyhf{}
\fancyhead[L]{\small\sffamily GPT 6 Astra Ultra research attempt}
\fancyhead[R]{\small Original catalogue IDs}
\fancyfoot[C]{\thepage}
\setlength{\parindent}{0pt}
\setlength{\parskip}{5pt plus 1pt}
\setlength{\emergencystretch}{3em}
\setcounter{tocdepth}{1}\setcounter{secnumdepth}{1}
\setlist[itemize]{leftmargin=3em,itemsep=4pt,topsep=4pt,parsep=0pt}
\allowdisplaybreaks
\newcommand{\N}{\mathbb{N}}
\newcommand{\Z}{\mathbb{Z}}
\newcommand{\Q}{\mathbb{Q}}
\newcommand{\R}{\mathbb{R}}
\newcommand{\C}{\mathbb{C}}
\newcommand{\F}{\mathbb{F}}
\newcommand{\A}{\mathbb{A}}
\newcommand{\PP}{\mathbb{P}}
\newcommand{\E}{\mathbb{E}}
\newcommand{\eps}{\varepsilon}
\newcommand{\dd}{\,\mathrm d}
\newcommand{\sref}[2]{\hyperlink{source:#1:#2}{[S#2]}}
\newif\ifshowcatalogue
\showcataloguetrue
\newcommand{\cataloguescope}[1]{\ifshowcatalogue
 \paragraph{Statement recorded in the supplied source.}
 #1\fi}
\begin{document}
% UnsolvedMath ID 1112; source catalogue code ALG-012

\setcounter{section}{1111}

\section{Existence of a Perfect Cuboid}\label{1112}

{\small\sffamily UnsolvedMath ID 1112 · Algebra}

\subsection{Definitions and mathematical statement}

\paragraph{Shared notation.}Unless a section says otherwise, $\N=\{1,2,\ldots\}$,
$\Z$ is the ring of integers, $\Q,\R,\C$ are the rational, real, and complex
fields, $[N]=\{1,\ldots,\lfloor N\rfloor\}$, and $\log$ is the natural
logarithm. For real functions of a parameter tending to infinity,
$f\ll g$ and $f=O(g)$ mean $|f|\leq Cg$ eventually for some fixed $C$;
$f\gg g$ reverses this comparison; $f=o(g)$ means $f/g\to0$;
$f\sim g$ means $f/g\to1$; and $f\asymp g$ means both order bounds.
A subscript on $O$ or $\ll$ specifies allowed constant dependence.
The expression $n^{o(1)}$ in an upper bound means growth smaller than
$n^\epsilon$ for each fixed $\epsilon>0$. Local definitions govern any
exception, finite-field convention, or use of the same symbol for another
object. An asymptotic claim is not automatically an effective algorithm.



Seek positive integers $a,b,c,u,v,w,t$ with
\[a^2+b^2=u^2,\quad a^2+c^2=v^2,\quad b^2+c^2=w^2,\quad a^2+b^2+c^2=t^2.\]
The three positive edge lengths rule out degenerate boxes. Rational solutions with positive edges are equivalent, by clearing denominators, to integer solutions. \sref{1112}{1} \sref{1112}{2}

\paragraph{Statement recorded in the supplied source.}
Does there exist a rectangular cuboid where all edges, face diagonals, and space diagonals have integer lengths? \sref{1112}{1}

\subsection{Short English statement}

Can a nondegenerate rectangular box have integer lengths for all its edges, all three face diagonals, and its space diagonal?

\subsection{Research attempt}

\paragraph{Scope and route.}
Every edge must be positive, and all four displayed square conditions
must hold simultaneously. Rational coordinates are equivalent after
scaling; a point with a zero edge is not a solution. The primary cuboid
paper [E1] explicitly distinguishes algebraic parametrizations from the
inverse rationality problems they leave. This attempt avoids treating an
algebraic parametrization as a rational solution, and instead studies
primitive congruences and a finite completion test for a fixed face.

\paragraph{Primitive reduction and parity.}
Divide all edges by $g=\gcd(a,b,c)$. Each diagonal is also divisible
by $g$: if its square is divisible by $g^2$, prime valuations imply
the diagonal is divisible by $g$. Thus it suffices to consider
$\gcd(a,b,c)=1$. Two odd edges would have squared sum two modulo four,
so at most one edge is odd. Primitivity excludes three even edges;
there is exactly one odd edge. Name it $a$.

The face diagonals involving $a$ are odd. Squares of odd integers are
one modulo eight. Therefore $a^2+b^2=u^2$ forces $b^2\equiv0\pmod8$,
which for an even $b$ means $4\mid b$. The same gives $4\mid c$.
These are stronger constraints than just ``two even edges.'' They are
still consistent: the Euler brick $(44,117,240)$ has one odd edge and
the other two divisible by four, yet its body diagonal is not integral.

\paragraph{Local obstructions at three and five.}
Squares modulo three are zero or one. If two edges were nonzero modulo
three, their face diagonal would have square two, impossible. Thus at
least two edges are divisible by three. Primitivity excludes all three,
so exactly two are divisible by three. This condition is independent
of which edge is odd.

At least one edge is divisible by five. Indeed if none were, their
square residues would all lie in $\{1,4\}$. Any two equal residues
sum to two or three modulo five, neither a square. With three residues
and only two choices, some pair is equal, contradicting its face
equation. Hence every primitive candidate has edge product divisible
by $2^4\cdot3^2\cdot5=720$; this product divisibility is a consequence
of the above separate assertions and not a claim that a single edge
contains every factor.

These modular arguments give necessary patterns rather than a local
contradiction. Passing all finitely many congruence tests does not
establish a rational point on the simultaneous equations.

\paragraph{A complete finite test for a fixed integral face.}
Fix positive $a,b$ with $a^2+b^2=u^2$. If a perfect completion by an
edge $c>0$ exists, its body diagonal satisfies
\[
 t^2-c^2=u^2,\qquad (t-c)(t+c)=u^2.
\]
Set $d=t-c$, $e=t+c$. Then $d,e$ are positive integers with
$de=u^2$, $d<e$, and $d\equiv e\pmod2$. Conversely every such factor
pair gives
\[
 c=(e-d)/2,\qquad t=(e+d)/2.
\]
The remaining two tests are exactly that $a^2+c^2$ and $b^2+c^2$
are squares. Thus all possible completions of a fixed integral face
can be tested by finitely many divisors of $u^2$, even without an a
priori bound on $c$. This is a completeness statement for that fixed
face. It does not give a finite list of faces to be considered.

For example the simplest face $(a,b,u)=(3,4,5)$ has only the positive
completion factor pair $(d,e)=(1,25)$, producing $c=12,t=13$.
But $3^2+12^2=153$ is not a square, so that face cannot occur in a
perfect cuboid. The other factor pair $(5,5)$ gives the forbidden edge
$c=0$. The argument eliminates all completions of this exact face,
not all rationally proportional faces: scaling a rational completion
may enlarge its integral face and introduce factor pairs not integral
at the original scale.

\paragraph{A rational parametrization and its remaining square tests.}
To see the rational issue directly, normalize $a=1$ and write
$x=b/a$, $y=c/a$. The two equations $1+x^2=U^2$ and
$1+y^2=V^2$ are parametrized by
\[
 x=\frac{r-r^{-1}}2,\quad U=\frac{r+r^{-1}}2,
 \qquad
 y=\frac{s-s^{-1}}2,\quad V=\frac{s+s^{-1}}2
\]
with positive rational $r,s>1$. This is complete because $r=U+x$ and
$s=V+y$. One still needs both $x^2+y^2=W^2$ and
$1+x^2+y^2=T^2$. Parametrizing $T^2-W^2=1$ with rational $v>1$
reduces these to
\[
 (r-r^{-1})^2+(s-s^{-1})^2=(v-v^{-1})^2.
\]
Conversely a rational triple $r,s,v>1$ satisfying this equation gives
all seven positive rational lengths. Thus it is an exact reduction,
but not a free three-parameter solution: the last equation is the
unresolved compatibility constraint. Taking a square root to solve for
$v$ over an algebraic extension does not prove that $v$ is rational.

\paragraph{Executed bounded search and a near miss.}
The script exhaustively checks sorted positive edges
$1\le a\le b\le c\le300$. It first requires $a^2+b^2$ to be a
square, then checks both other faces and finally the body diagonal.
The only Euler bricks in this box are $(44,117,240)$ and
$(240,252,275)$; neither has an integral body diagonal. The first has
face diagonals $125,244,267$ and squared body diagonal $73225$,
strictly between $270^2$ and $271^2$. All these tests use integer
square roots. Files: \texttt{checks.py}, \texttt{results.json} under
\path{_Local/Erdos_problems/UnsolvedMath/attempt_audit_2026_10_01/number_theory/}.
The code executes the bounded edge search, not an unbounded rational
point algorithm or the fixed-face factor-pair strategy for every face.

\paragraph{Remaining gap and outcome.}
The primitive congruences, finite fixed-face completion criterion, and
rational compatibility equation are proved. None excludes every
positive rational solution of the final equation, and no such solution
was found. A finite height bound would be a major additional ingredient;
it is not supplied by clearing denominators. \textbf{Outcome: UNRESOLVED.}
The near misses satisfy only proper subsets of the required equations.

\subsection{Sources}

{\small

\begin{itemize}

\item[\hypertarget{source:1112:1}{S1}] ULAM.AI, \emph{UnsolvedMath}, supplied \texttt{problems.json}, record 1112 (ALG-012), statement and background. The embedded literature-review date is 2026-08-17; that is the archive’s review date, not a fresh certification. Dataset landing page: \url{https://huggingface.co/datasets/ulamai/UnsolvedMath}.

\item[\hypertarget{source:1112:2}{S2}] Perfect cuboid overview (). \url{https://en.wikipedia.org/wiki/Perfect_cuboid}. Bibliographic information imported from the supplied record.


\item[E1] John Ramsden and Ruslan Sharipov,
\emph{On two algebraic parametrizations for rational solutions of the
cuboid equations}, arXiv:1208.2587 (2012).
\url{https://arxiv.org/abs/1208.2587}. Abstract inspected 1 October 2026.
Its inverse-problem distinction is used for scope; its full constructions
are not premises of the elementary arguments here.

\end{itemize}

}

\label{end:1112}
\end{document}
